% Ad M. Brutum 1.4 (ad-brutum-01-04)
% Source: https://raw.githubusercontent.com/PerseusDL/canonical-latinLit/master/data/phi0474/phi059/phi0474.phi059.perseus-lat1.xml
% Fetched by scripts/fetch_latin.py

% Dateline (Perseus): Scr. Dyrrachi iii antprid. id. Mai. a. 711 (43) .

\ciceroLetterOpener{Brutus Ciceroni salutem}

\ciceroSection{1} quanta sim laetitia adfectus cognitis rebus Bruti nostri et consulum facilius est tibi existimare quam mihi scribere. Cum alia laudo et gaudeo accidisse, tum quod Bruti eruptio non solum ipsi salutaris fuit sed etiam maximo ad victoriam adiumento.

\ciceroSection{2} quod scribis mihi trium Antoniorum unam atque eandem causam esse, quid ego sentiam mei iudici esse: statuo nihil nisi hoc, senatus aut populi Romani iudicium esse de iis civibus qui pugnantes non interierint. at hoc ipsum inquies inique facis qui hostilis animi in rem publicam homines civis appelles. immo iustissime. quod enim nondum senatus censuit nec populus Romanus iussit, id adroganter non praeiudico neque revoco ad arbitrium meum. illud quidem non muto, quod ei quem me occidere res non coegit neque crudeliter quicquam eripui neque dissolute quicquam remisi habuique in mea potestate quoad bellum fuit. multo equidem honestius iudico magisque quod concedere possit res publica miscrorum fortunam non inscctari quam infinite tribucre potentibus quac cupiditatem et adrogantiam incendere possint.

\ciceroSection{3} qua in re, Cicero, vir optime atque fortissime mihique merito et meo nomine et rei publicae carissime, nimis credere videris spei tuae statimque, ut quisque aliquid recte fecerit, omnia dare ac permittere, quasi non liceat traduci ad mala consilia corruptum largitionibus animum. quae tua est humanitas, aequo animo te moneri patieris, praesertim de communi salute; facies tamen quod tibi visum fuerit; etiam ego, cum me docueris * *
